\section{The category of multiplets} \label{sec: multiplets}
In~\cite{perspectives}, the notion of a \emph{$\fg$-multiplet} for a super Lie algebra $\fg$ was formalized. Let us quickly review the relevant definitions, and then move on to define morphisms between $\fg$-multiplets and study the resulting category $\Mult_{\fg}$. We start with the case of strict multiplets and strict morphisms and then discuss the homotopy-theoretic generalization.

In brief, by a \emph{multiplet} on a smooth manifold $X$ we will mean a graded vector bundle $E$ on $X$ together with a differential $D$ and an action $\rho$ of an appropriate super Lie algebra $\fg$, compatible with an action of the even part of $\fg$ on the base manifold $X$. The reader should consider the prototypical situation where $X = \RR^n$, and $\fg$ is a super Lie algebra whose even part is the Poincar\'e algebra of infinitesimal isometries of $\RR^n$. A multiplet can be thought of as modelling the perturbations of a classical field in the BV formalism.

\begin{rmk}
	We would like to clearly emphasize here the distinction between multiplets and theories in this context. In general, the data of a multiplet is \emph{less data} than the data of a classical field theory: it does not include enough data to specify dynamics. The data of a perturbative classical supersymmetric field theory in the BV formalism includes the data of a multiplet, together with some additional data. If a multiplet is equipped with the additional data of a~\emph{BV bracket} (an odd shifted symplectic structure), it can be viewed as modelling a \emph{free} theory in the BV formalism. Such structures can easily be constructed from multiplets by taking the sum of a multiplet with a shift of its dual.	If one wishes to include non-trivial interactions, one must additionally equip the sections of the bundle $E$ with an $L_\infty$ structure, compatibly with the supersymmetry action.
	
	So, while the results of this paper do apply to BV field theories, when we speak of \emph{multiplets} we are referring to a less constrained notion, and to obtain a genuine field theory from a multiplet one needs to specify strictly more data.
\end{rmk}

\subsection{Gradings}
\label{ssec:gradings}
Many objects appearing throughout this work (e.g., vector spaces, vector bundles, Lie algebras etc.) will carry a grading by $\ZZ\times \ZZ/2\ZZ$ as well as a differential of bidegree $(1,+)$. We refer to such objects with the prefix ``dgs'' standing for ``differential graded super''. Typically these gradings have a clear physical meaning: the integer grading refers to the ghost number or cohomological degree, while the $\ZZ/2\ZZ$-grading corresponds to the intrinsic parity (fermion number modulo two). The total parity denotes the $\ZZ/2\ZZ$-grading which arises by forgetting the $\ZZ$-grading to~$\ZZ/2\ZZ$ and then totalizing with the intrinsic parity. It is the total parity which governs all signs.

It will often be convenient to introduce an additional piece of data: an auxiliary action of a one-dimensional abelian Lie algebra $\RR$ on a dgs vector space. Many of our examples will be built from the following simple observation.

\begin{eg}
	Let $\fn$ be a dg Lie algebra concentrated in degrees 1 and 2. For all non-negative integers $k$, there is a natural action of $\RR$ on $\Sym^k(\fn^*)$, where $\RR$ acts on duals to the degree 1 and 2 summands of $\fn^*$ with weight $-1$ and $-2$ respectively.
\end{eg}

The results in this paper will make sense for arbitrary $\RR$-equivariant dgs vector spaces, but our most common examples will have the following behavior.

\begin{dfn}
	A \emph{lifted} dgs vector space is an $\RR$-equivariant dgs vector space where the $\RR$-action has integral weights, and the $\ZZ/2\ZZ$-grading coincides with the $\RR$-weight modulo two.
\end{dfn}

Let us now introduce some notation that we will use when we discuss lifted dgs vector bundles on a smooth manifold $X$. Let us write
\begin{equation*}
E = \bigoplus_{(w,d) \in \ZZ^2} E^{w,d}
\end{equation*}
for a lifted dgs vector bundle, where the first index indicates the decomposition of $E$ by $\RR$-weight, and the second index indicates internal $\ZZ$-degree on $E$. In some of our calculations, particular in~Section~\ref{sec: examples}, we will place the summands in a two-dimensional array, where the two coordinates are given by $w-d$ and $d$. For example,
\begin{equation*}
E =
\begin{bmatrix}
\cdots &E^{-1,0} &E^{0,0} &E^{1,0} &E^{2,0}&\cdots\\
\cdots &E^{0,1} &E^{1,1} &E^{2,1} &E^{3,1}&\cdots\\
\cdots &E^{1,2} &E^{2,2} &E^{3,2} &E^{4,2} &\cdots
\end{bmatrix}.
\end{equation*}
So we recover $w$ as the total degree with respect to this sheared bigrading. The overall parity is then determined by the column.

If the differential on our dgs vector bundle decomposes as a sum of homogeneous differential operators of order $k$, say $D = \sum_{k \ge 0} D_k$, the summand $D_k$ acts---with respect to the sheared grading---with bidegree $(2k-1, 1)$. So the homogeneous summands all increase the vertical degree by 1, but they may modify the horizontal degree by any odd integer $\ge -1$.

\subsection{Strict multiplets}

For a dgs vector bundle\footnote{Note that each graded piece $E^k$ is a finite rank vector bundle, but the total rank of $E$ may still be infinite.} $(E,D)$ over a base smooth manifold $X$, we denote the space of global smooth sections by $\cE = \Gamma(X,E)$. The endomorphisms $\End(\cE)$ form a dgs Lie algebra, where the bracket is given by the commutator and the differential is $[D,-]$. Inside $(\End(\cE), [D,-])$ there is a sub dgs Lie algebra consisting of all endomorphisms which act on sections via differential operators. We denote this subalgebra by $(\cD(E), [D,-])$.
\begin{dfn}
	A \emph{strict local dgs $\fg$-module} is a triple $(E,D,\rho)$ where $(E,D)$ is a dgs vector bundle and
	\begin{equation*}
	\rho \colon \ \fg \longrightarrow (\cD(E) , [D,-] )
	\end{equation*}
	is a map of dgs Lie algebras. Here the super Lie algebra $\fg$ is viewed as a dgs Lie algebra in cohomological degree zero with trivial differential.
\end{dfn}
\begin{rmk}
	This definition (as well as many of the following) has a natural generalization for~$\fg$ a dgs Lie algebra. Since we are ultimately interested in the pure spinor superfield formalism, we restrict our attention to super Lie algebras with no cohomological grading.
\end{rmk}
Note that, since a super Lie algebra $\fg$ has no differential, $\rho$ commutes with the differential on the dgs vector bundle,
\begin{equation*}
[D,\rho(x)] = 0 \qquad \forall x \in \fg .
\end{equation*}
It is standard to encode a $\fg$-module structure $\rho$ on $(E,D)$ as a dgs Lie algebra structure on the direct sum $\fg \oplus \cE$. Concretely we set for the unary and binary operations
\begin{align}
&\mu_1(x_1,\sigma_1) = (0, D\sigma_1),\nonumber \\
&\mu_2((x_1,\sigma_1),(x_2,\sigma_2)) = \big([x_1,x_2] , \rho(x_1)\sigma_2 - (-1)^{|\sigma_1| |x_2|} \rho(x_2)\sigma_1 \big) ,\label{eq: lie structure on sum}
\end{align}
where $x_1,x_2 \in \fg$ and $\sigma_1,\sigma_2 \in \cE$.

There is an obvious notion of morphisms between strict local $\fg$-modules.
\begin{dfn}
	A \emph{strict morphism} of strict local $\fg$-modules	$(E,D,\rho)$ and $(E',D',\rho')$ is a map of cochain complexes
	\begin{equation*}
	\psi \colon \ \cE \longrightarrow \cE'
	\end{equation*}
	realized by differential operators such that
	\begin{equation*}
	\psi \circ \rho(x) = \rho'(x) \circ \psi
	\end{equation*}
	for all $x \in \fg$.
\end{dfn}
A strict morphism $\psi \colon (E,D,\rho) \longrightarrow (E',D',\rho')$ gives rise to a strict morphism of the associated dgs Lie algebras by setting $\tilde{\psi} = \id_{\fg} \times \psi \colon \fg \oplus \cE \longrightarrow \fg \oplus \cE'$. Conversely it is easy to check that every strict morphism of dgs Lie algebras of that form gives rise to a strict morphism of $\fg$-modules. We call $\psi$ a quasi-isomorphism if it is a quasi-isomorphism of dgs vector bundles; equivalently $\tilde{\psi}$ is a quasi-isomorphism of dgs Lie algebras.

Now, suppose that $(E,D)$ is a dgs vector bundle over an affine space $V$ (so, in the notation used above, $V=X$). Let us additionally assume that $\fg$ is equipped with a map of super Lie algebras
\begin{equation*}
\phi \colon \ \aff(V) \longrightarrow \fg .
\end{equation*}
In essence, when we refer to strict $\fg$-\emph{multiplets}, we are referring to strict local $\fg$-modules for which the affine transformations acts in a geometric way.
\begin{dfn}
	A \emph{strict $\fg$-multiplet} $(E,D,\rho)$ on $V$ is an affine dgs vector bundle over $V$ equipped with a strict local $\fg$-module structure
	\begin{equation*}
	\rho \colon \ \fg \longrightarrow \cD(E)
	\end{equation*}
	such that the pullback of the module structure along~$\phi$ agrees with the natural action on sections of the affine vector bundle. Concretely, this means that the following diagram commutes:
	\begin{equation} \label{eq: geom action}
	\begin{tikzcd}
	\fg \arrow[r, "\rho"] & \cD(E). \\
	\aff(V) \arrow[ru, "\mathrm{aff}" '] \arrow[u,"\phi"]
	\end{tikzcd}
	\end{equation}
\end{dfn}
Here $\mathrm{aff}$ denotes the natural action of the affine algebra on $E$.
We define the category of strict multiplets with strict morphisms to be the full subcategory of the category of strict local $\fg$-modules with objects strict $\fg$-multiplets. We denote this category by $\Mult_{\fg}^{\mathrm{strict}}$.
%
\begin{rmk}
	The condition~\eqref{eq: geom action} stating that the translations act geometrically heavily depends on $\phi$. In particular, since the action $\mathrm{aff}$ is effective, the corresponding category of $\fg$-multiplets is trivial if $\phi$ is not injective.
\end{rmk}
\begin{rmk}
	Note that a morphism of strict local $\fg$-modules is automatically compatible with the action of $\aff(V)$. For example let $\psi \colon (E,D,\rho) \longrightarrow (E',D',\rho')$ be a strict morphism. Then we have
	\begin{equation*}
	\psi \circ \mathrm{aff}(x) = \psi \circ \rho(\phi(x)) = \rho'(\phi'(x)) \circ \psi = \mathrm{aff}'(x) \circ \psi .
	\end{equation*}
	Therefore it is sensible to define $\Mult_\fg^{\text{strict}}$ as a full subcategory of strict local $\fg$-modules.
\end{rmk}
%
\subsection{Homotopy theory of multiplets}
%
It is well known that various strict algebraic structures, such as associative algebras or Lie algebras, admit homotopy-theoretic generalizations ($A_\infty$ or $L_\infty$ algebras in these cases). The same is true for strict $\fg$-module structures,
which generalize to $L_\infty$ $\fg$-modules. In the context of this work, we will refer to such homotopy module structures simply as module structures and choose to emphasize whenever a module is strict instead. In this spirit we can easily define (not necessarily strict) local $\fg$-modules and $\fg$-multiplets by replacing Lie maps in the above definitions with $L_\infty$ maps.

Thus, a local $\fg$-module is just a dgs vector bundle $(E,D)$ together with an $L_\infty$ map of $L_\infty$ algebras
\begin{equation*}
\rho \colon \ \fg \rightsquigarrow \cD(E) .
\end{equation*}
Recall that this means there are component maps
\begin{equation*}
\rho^{(k)} \colon \ \fg^{\otimes k} \longrightarrow \cD(E)[1-k], \qquad k \geq 1
\end{equation*}
satisfying a series of compatibility relations the lowest of which reads
\begin{equation*}% \label{eq: L-action}
\big[\rho^{(1)}(x),\rho^{(1)}(y)\big] - \rho^{(1)}([x,y]) =\big[D, \rho^{(2)}(x,y)\big] \qquad \forall x,y \in \fg.
\end{equation*}
Clearly a local $\fg$-module is strict if and only if $\rho^{(k)} = 0$ for all $k\geq2$.

\begin{rmk}
	As we discussed in the introduction, considering this homotopical version of a~module for a super Lie algebra is very natural in supersymmetric field theory. In most supersymmetric classical field theories, when one describes the most natural multiplet of fields, there is only a strict action of the supersymmetry algebra on-shell, i.e., after imposing the equations of motion. If one, however, allows $L_\infty$ actions, it is possible to describe an $L_\infty$ action of the supersymmetry algebra on the space of fields in the BV formalism directly, without introducing auxiliary fields. This idea goes back to the beginnings of the BV formalism; applications to global symmetries are discussed, for example, in~\cite{BaulieuSusy}. The example of super Yang--Mills theory is discussed rigorously and systematically in all dimensions in~\cite{ESW}.
\end{rmk}


Similarly to the strict case, we can conveniently encode a $\fg$-module structure $\rho$ as an~$L_\infty$ structure on $\fg \oplus \cE$. To this end we supplement the operations~\eqref{eq: lie structure on sum} by the following brackets for $k\geq3$ (for details we refer to~\cite{Allocca, Lada})
\begin{equation*}
\mu_k((x_1,v_1),\dots,(x_k,v_k)) = \left( 0 , \sum_{i=1}^{k} \pm \rho^{(k-1)}\big(x_1,\dots,\hat{x}_i,\dots,x_k\big)v_i \right).
\end{equation*}
One can define a morphism $\psi \colon (E,D,\rho) \longrightarrow (E',D',\rho')$ of local $\fg$-modules by component maps
\begin{equation*}
\psi_n \colon \ \fg^{\otimes n-1} \otimes \cE \longrightarrow \cE' ,
\end{equation*}
which are given by differential operators and satisfy a series of compatibility relations (see~\cite{Allocca} for details). We recover strict morphisms by restricting to those where the only component map is $\psi_1$. Again, we can describe such morphisms by morphisms of the associated $L_\infty$ algebras
\begin{equation*}
\tilde{\psi} \colon \ \fg \oplus \cE \rightsquigarrow \fg \oplus \cE'
\end{equation*}
by supplementing $\tilde{\psi}_1 = \id_{\fg} \times \psi_1$ with
\begin{equation*}
\tilde{\psi}_k ((x_1,\sigma_1) , \dots , (x_k, \sigma_k)) = \left(0, \sum_{i=1}^{k} \pm \psi_k\big(x_1, \dots , \hat{x}_i , \dots , x_k , \sigma_i \big)\right)
\end{equation*}
for $k \geq 2$. $\psi$ is a quasi-isomorphism of local $\fg$-modules if $\psi_1$ is a quasi-isomorphism of cochain complexes, or equivalently if $\tilde{\psi}$ is a quasi-isomorphism of $L_\infty$ algebras. In general, encoding module structures as $L_\infty$ structures is very convenient because it allows the use of many known tools from the theory of $L_\infty$ algebras, like homotopy transfer for instance.

\begin{rmk}
	An equivalent realization of the notion of a non-strict morphism between a pair of local $L_\infty$ algebras $\fg,\fg'$ is provided by considering the Chevalley--Eilenberg chain complex~$C_\bullet(\fg)$ of an $L_\infty$ algebra $\fg$. This is a cocommutative dg coalgebra whose underlying graded coalgebra is $\Sym^\bullet(\fg[-1])$, with differential induced from the $L_\infty$ structure on $\fg$ as a sum of terms given by the $L_\infty$ brackets on $\fg$. The data of a morphism $\psi \colon \fg \to \fg'$ of local $L_\infty$ algebras is equivalent to that of a morphism of cocommutative dg coalgebras
	\[\psi' \colon \ C_\bullet(\fg) \to C_\bullet(\fg')\]
	given by differential operators. This interpretation allows us to view local $L_\infty$ algebras as objects of a dg-category, so that we can discuss (for example) homotopies between $L_\infty$ algebra morphisms. When the source $\fg$ is finite-dimensional, we can dually consider morphisms of commutative dg algebras $C^\bullet(\fg') \to C^\bullet(\fg)$. This is for example the case for the local $\fg$-module structures $\rho \colon \fg \rightsquigarrow \cD(E)$ appearing in the definition of multiplets.
\end{rmk}

Let us now give the definition of a (not necessarily strict) multiplet.
\begin{dfn}
	A \emph{$\fg$-multiplet} $(E,D,\rho)$ is an affine dgs vector bundle over $V$ equipped with a local $\fg$-module structure
	\begin{equation*}
	\rho\colon \ \fg \rightsquigarrow \cD(E)
	\end{equation*}
	such that the pullback of the module structure along~$\phi$ agrees strictly with the natural action on sections of the affine vector bundle. Concretely, this means that the following diagram commutes:
	\begin{equation*}
	\begin{tikzcd}
	\fg \arrow[r, "\rho^{(1)}"] & \cD(E). \\
	\aff(V) \arrow[ru, "\mathrm{aff}" '] \arrow[u,"\phi"]
	\end{tikzcd}
	\end{equation*}
	We define the dg-category of $\fg$-multiplets to be the full subcategory of local $\fg$-modules with objects being $\fg$-multiplets and denote it by $\Mult_{\fg}$.
\end{dfn}

\begin{rmk}
	Note that, in particular, the action $\rho \circ \phi$ obtained by restricting the $L_\infty$ action to affine transformations is necessarily strict.
\end{rmk}

We will sometimes wish to refer to the full subcategory of strict multiplets, but allowing arbitrary morphisms.
\begin{dfn}
	Denote by $\Mult_{\fg}^{\text{\rm strict-ob}}$ the full sub dg-category of $\Mult_{\fg}$ generated by strict multiplets.
\end{dfn}


We can obtain homotopy categories from the dg-categories $\Mult_{\fg}$ as well as $\Mult_{\fg}^{\mathrm{strict}}$ by replacing the hom space $\Hom_{\Mult_{\fg}}(E, E')$ by its zeroth cohomology $\mathrm H^0(\mathrm{Hom}_{\Mult_{\fg}}(E, E'))$. We denoted the resulting categories by $\mathrm{Ho}(\Mult_{\fg})$ and $\mathrm{Ho}(\Mult_{\fg}^{\mathrm{strict}})$. Speaking physically, quasi-isomorphisms of $\fg$-multiplets correspond to \emph{perturbative} equivalences of multiplets, where by ``perturbative'' here we mean equivalences of the derived formal neighborhoods of a point in the classical moduli field space. Therefore isomorphism classes in the homotopy category correspond to perturbatively distinct multiplets.


\begin{rmk}
	Because every $L_\infty$ algebra can be strictified \cite[Part II, Corollary 1.6]{KrizMay}, the natural inclusion
	\begin{equation*}
	\mathrm{Ho}\big(\Mult_{\fg}^{\text{\rm strict-ob}}\big) \to \mathrm{Ho}(\Mult_{\fg})
	\end{equation*}
	is an equivalence of categories, where $\Mult_{\fg}^{\text{\rm strict-ob}}$ denotes the full subcategory of multiplets spanned by strict objects.
\end{rmk}
